\chapter{Coordenadas angulares conjugadas y descomposición espectral}
\label{chap:canales}

El paso desde el cierre aritmético de \(\alpha\) hacia la jerarquía
dimensional se realiza mediante dos coordenadas angulares. La primera procede
directamente de \(\alphaHMT\); la segunda es la fase orientada producida por
la realización analítica regional de la microfibra quíntuple. Ambas llegan a
Mecánica Dimensional ya construidas, antes de introducir la carta exponencial
que generará los momentos globales.

\section{Transferencia de las coordenadas angulares derivadas en HMT}

Las Partes I--III construyeron dependencias distintas con raíz común en
\(\alpha_{\rm HMT}\). El carácter pre-retorno procede de
\[
(\alpha_{\rm HMT},\varphi;\,9,5,4,12,54,\text{ regla graduada})
\longrightarrow H_5,
\]
mientras que la coordenada angular y el determinante proceden de
\[
\alpha_{\rm HMT}\longrightarrow
A=1000\alpha_{\rm HMT}\,{}^\circ,
\qquad
(\alpha_{\rm HMT},\pi)\longrightarrow
D_A=e^{-100\pi\alpha_{\rm HMT}/9}.
\]
La microfibra regional fija, mediante una dependencia distinta,
\(\mathcal F_\pi\to\rho_\pi=169\), y con ella
\[
\mathcal C_\pi(\alpha_{\rm HMT})
=169\alpha_{\rm HMT}^{6}
+\frac{D_A\alpha_{\rm HMT}^{7}}
       {1-D_A\alpha_{\rm HMT}}.
\]
La rama de \(H_5\) y la corrección regional confluyen sólo en el paso
siguiente:
\[
\begin{aligned}
y_*&=\frac{\pi}{90}(H_5+\alpha_{\rm HMT})
     -\mathcal C_\pi(\alpha_{\rm HMT}),\\
C^*&=\frac{180}{\pi}y_*.
\end{aligned}
\]
En adelante, \(A\) denota la coordenada angular directa y \(C^*\) la
coordenada angular regional conjugada; estas denominaciones fijan sus tipos
sin introducir una segunda regla de selección.
Por tanto, \(H_5\) no es una salida de \(A\). Aquí \(y_*\) es la coordenada
analítica en radianes y \(C^*\) su coordenada sexagesimal. Mecánica
Dimensional recibe estas salidas ya determinadas; no vuelve a definirlas.
La coordenada \(A\) no es entrada directa de la fórmula de \(C^*\): ambos
objetos se acoplan por primera vez en \((A,C^*)\mapsto q_\pm\). La igualdad
\(\alpha_{\rm HMT}=A/1000\) es únicamente el cambio de coordenada inverso.
En esta Parte conservamos las notaciones
\[
\begin{aligned}
 A&=7.297352569283800997285105472380663^\circ,\\
 \Cstar&=2.123738338968645724261951380393\ldots{}^\circ.
\end{aligned}
\]
El redondeo empleado en las tablas es
\(2.123738338968646^\circ\). El entero \(169\), el determinante \(D_A\) y
la recurrencia regional pertenecen a la demostración anterior; ninguna
comparación SI ni el funcional hexada--octada intervienen en su generación.

La
identidad angular
\[
 \etaRet=\frac{\Cstar}{2}-\alphaAngle,
 \qquad
 \alphaAngle:=\alphaHMT\,{\rm deg},
\]
define la coordenada angular adimensional posterior al retorno; no se emplea
para generar \(\Cstar\). Equivalentemente,
una vez efectuada la construcción regional,
\[
 \Cstar=2(\etaRet+\alphaAngle).
\]
Esta última igualdad es la forma inversa de la misma relación y no una
segunda regla de selección de \(C^*\).

La realización regional corrige \(H_5\) antes de producir \(\etaRet\); por
tanto, \(H_5\),
\(\etaRet\) y \(\Cstar\) tienen tipos matemáticos distintos. El símbolo
\(\hbarHMT\) se reserva para la magnitud sobre la recta de acción obtenida
sólo después de aplicar la escala \(10^{-34}\mathcal U_S\). La comparación
con una mantisa del SI pertenece al control metrológico externo y deja
invariante la construcción anterior.

\section{Fases internas y carta exponencial}

Las fases abstractas asociadas a los dos ángulos son
\[
 a=\frac A{360},\qquad c=\frac{\Cstar}{360}
 \quad\text{en }\RR/\ZZ.
\]
Antes de aplicar una exponencial real se eligen los representantes canónicos
\(\widetilde a,\widetilde c\in[0,1)\). La exponencial actúa sobre esos
representantes, no directamente sobre las clases del cociente.
La conversión a radianes se introduce únicamente al elegir la carta
arquimediana
\[
 x:=A_{\rm rad}=\frac{\pi A}{180},\qquad
 y:=\Cstarrad=\frac{\pi\Cstar}{180},
\]
\[
 q_-=e^{-x+y},\qquad q_+=e^{-x-y}.
\]
El factor \(\pi/180\) pertenece a este cambio analítico de coordenadas. El
cierre aritmético de \(\alpha\), formulado antes de esta conversión, es
independiente de dicha elección de carta angular.

\section{Álgebra involutiva central}

Consideremos la representación regular bidimensional del álgebra real
generada por una involución. En una base espectral fijamos
\[
 R=\operatorname{diag}(1,-1),\qquad I=I_2,
\]
de modo que ambos subespacios propios tienen multiplicidad uno. Sean
\[
 P_+=\frac{I+R}{2},\qquad P_-=\frac{I-R}{2}
\]
sus proyectores espectrales. El elemento central determinado por las
coordenadas angulares es
\[
 B_{A,\Cstar}=AI+\Cstar R
 =(A+\Cstar)P_++(A-\Cstar)P_-.
\]
La involución \(\Cstar\mapsto-\Cstar\) conserva los proyectores fijos
\(P_\pm\) e intercambia los dos coeficientes espectrales ---equivalentemente,
los canales asociados---. Al aplicar el exponencial se obtiene
\[
 T=e^{-xI-yR}=q_+P_++q_-P_-.
\]
Por tanto, los dos canales son los autovalores de un único elemento del
álgebra generada por \(I\) y \(R\).

\section{Invertibilidad de la carta de canales}

\begin{theorem}[Reconstrucción de las coordenadas angulares]
Para \(q_\pm>0\),
\[
 x=-\frac12\log(q_+q_-),
 \qquad
 y=\frac12\log\frac{q_-}{q_+}.
\]
En consecuencia, la aplicación
\((A,\Cstar)\mapsto(q_+,q_-)\) es inyectiva. El intercambio
\(q_+\leftrightarrow q_-\) conserva \(A\) y transforma
\(\Cstar\mapsto-\Cstar\).
\end{theorem}
\begin{proof}
La multiplicación y el cociente de las dos exponenciales dan
\(q_+q_-=e^{-2x}\) y \(q_-/q_+=e^{2y}\). Al tomar logaritmos se recuperan
\(x\) e \(y\). Como \(\pi/180\neq0\), quedan determinados ambos ángulos.
\end{proof}

\section{Identidades de adición y recurrencia}

En esta representación bidimensional, los momentos par e impar del operador
\(T\) son
\[
 c_n=\Tr(T^n)=q_-^n+q_+^n,
 \qquad
 s_n=-\Tr(RT^n)=q_-^n-q_+^n.
\]

La hipótesis de multiplicidad uno forma parte del dominio de estas fórmulas.
En una representación con multiplicidades \(m_+\) y \(m_-\), las trazas
correspondientes serían
\(m_+q_+^n+m_-q_-^n\) y
\(m_-q_-^n-m_+q_+^n\), respectivamente.

\begin{theorem}[Leyes de composición de los momentos]
Para \(m,n\geq0\),
\[
 c_{m+n}=\frac{c_mc_n+s_ms_n}{2},
 \qquad
 s_{m+n}=\frac{s_mc_n+c_ms_n}{2},
\]
\[
 c_n^2-s_n^2=4e^{-2nx},
\]
y ambas sucesiones satisfacen
\[
 X_{n+2}=c_1X_{n+1}-e^{-2x}X_n.
\]
Además,
\[
 \partial_Ac_n=-\frac{n\pi}{180}c_n,
 \quad
 \partial_As_n=-\frac{n\pi}{180}s_n,
\]
\[
 \partial_{\Cstar}c_n=\frac{n\pi}{180}s_n,
 \quad
 \partial_{\Cstar}s_n=\frac{n\pi}{180}c_n.
\]
\end{theorem}
\begin{proof}
Se expanden los productos de
\(q_-^m\pm q_+^m\) con \(q_-^n\pm q_+^n\). La identidad cuadrática resulta
de \((u+v)^2-(u-v)^2=4uv\). La recurrencia corresponde al polinomio
característico \(z^2-c_1z+q_-q_+\), y las derivadas se obtienen mediante la
regla de la cadena.
\end{proof}

\begin{corollary}[Determinación de la jerarquía completa]
Fijados \(A\) y \(\Cstar\), todas las parejas \((c_n,s_n)\) quedan
determinadas por la recurrencia. Ninguna especie física introduce un canal
adicional.
\end{corollary}

Finalmente,
\[
 s_n=2e^{-nx}\sinh(ny),
 \qquad
 c_n=2e^{-nx}\cosh(ny).
\]
La coordenada \(A\) determina el decaimiento común y \(\Cstar\) la separación
orientada. El orden decimal de la recta de acción no se introduce en estas
fórmulas. Se deduce en el \cref{chap:escala-accion} a partir del conúcleo
integral del bloque local de Hadamard. La elección posterior de una base
expresada en \(\mathrm{J\,s}\) pertenece al transductor metrológico y no
altera las identidades adimensionales.
